% to have line numbers
%\RequirePackage{lineno}
\documentclass[10pt, letterpaper]{article}      
\usepackage[margin=.1cm,font=small,labelfont=bf]{caption}[2007/03/09]
%\usepackage{endnotes}
%\let\footnote=\endnote
\usepackage{setspace}
\usepackage{longtable}                        
\usepackage{anysize}                          
\usepackage{natbib}                           
%\bibpunct{(}{)}{,}{a}{,}{,}                   
\bibpunct{(}{)}{,}{a}{}{,}                   
\usepackage{amsmath}
\usepackage[% draft,
pdftex]{graphicx} %draft is a way to exclude figures                
\usepackage{epstopdf}
\usepackage{hyperref}                             % For creating hyperlinks in cross references
\hypersetup{pdfborder={0 0 0.4}} %have nice light boxes on refs

% \usepackage[margins]{trackchanges}

% \note[editor]{The note}
% \annote[editor]{Text to annotate}{The note}
%    \add[editor]{Text to add}
% \remove[editor]{Text to remove}
% \change[editor]{Text to remove}{Text to add}

%TODO make it more standard before submission: \marginsize{2cm}{2cm}{1cm}{1cm}
\marginsize{1cm}{1cm}{.5cm}{.5cm}%{left}{right}{top}{bottom}   
					          % Helps LaTeX put figures where YOU want
 \renewcommand{\topfraction}{1}	                  % 90% of page top can be a float
 \renewcommand{\bottomfraction}{1}	          % 90% of page bottom can be a float
 \renewcommand{\textfraction}{0.0}	          % only 10% of page must to be text

 \usepackage{float}                               %latex will not complain to include float after float

\usepackage[table]{xcolor}                        %for table shading
\definecolor{gray90}{gray}{0.90}
\definecolor{orange}{RGB}{255,128,0}

\renewcommand\arraystretch{.9}                    %for spacing of arrays like tabular

%-------------------- my commands -----------------------------------------
\newenvironment{ig}[1]{
\begin{center}
 %\includegraphics[height=5.0in]{#1} 
 \includegraphics[height=3.3in]{#1} 
\end{center}}

 \newcommand{\cc}[1]{
\hspace{-.13in}$\bullet$\marginpar{\begin{spacing}{.6}\begin{footnotesize}\color{blue}{#1}\end{footnotesize}\end{spacing}}
\hspace{-.13in} }

%-------------------- END my commands -----------------------------------------



%-------------------- extra options -----------------------------------------

%%%%%%%%%%%%%
% footnotes %
%%%%%%%%%%%%%

%\long\def\symbolfootnote[#1]#2{\begingroup% %these can be used to make footnote  nonnumeric asterick, dagger etc
%\def\thefootnote{\fnsymbol{footnote}}\footnote[#1]{#2}\endgroup}	%see: http://help-csli.stanford.edu/tex/latex-footnotes.shtml

%%%%%%%%%%%
% spacing %
%%%%%%%%%%%

% \abovecaptionskip: space above caption
% \belowcaptionskip: space below caption
%\oddsidemargin 0cm
%\evensidemargin 0cm

%%%%%%%%%
% style %
%%%%%%%%%

%\pagestyle{myheadings}         % Option to put page headers
                               % Needed \documentclass[a4paper,twoside]{article}
%\markboth{{\small\it Politics and Life Satisfaction }}
%{{\small\it Adam Okulicz-Kozaryn} }

%\headsep 1.5cm
% \pagestyle{empty}			% no page numbers
% \parindent  15.mm			% indent paragraph by this much
% \parskip     2.mm			% space between paragraphs
% \mathindent 20.mm			% indent math equations by this much

%%%%%%%%%%%%%%%%%%
% extra packages %
%%%%%%%%%%%%%%%%%%

\usepackage{datetime}


\usepackage[latin1]{inputenc}
\usepackage{tikz}
\usetikzlibrary{shapes,arrows,backgrounds}
\usetikzlibrary{trees}



%\usepackage{color}					% For creating coloured text and background
%\usepackage{float}
\usepackage{subfig}                                     % for combined figures

\renewcommand{\ss}[1]{{\colorbox{blue}{\bf \color{white}{#1}}}}
\newcommand{\ee}[1]{\endnote{\vspace{-.10in}\begin{spacing}{1.0}{\normalsize #1}\end{spacing}\vspace{.20in}}}
\newcommand{\emd}[1]{\ExecuteMetaData[/tmp/tex]{#1}} % grab numbers  from stata

%TODO before submitting comment this out to get 'regular fornt'
\usepackage{sectsty}
\allsectionsfont{\normalfont\sffamily}
\usepackage{sectsty}
\allsectionsfont{\normalfont\sffamily}
\renewcommand\familydefault{\sfdefault}

%\usepackage[margins]{trackchanges}
\usepackage{rotating}
\usepackage{catchfilebetweentags}

\usepackage{abstract}
\renewcommand{\abstractname}{}    % clear the title
\renewcommand{\absnamepos}{empty} % originally center


\usepackage{pgfgantt}


% -------------------- END extra options -----------------------------------------
\date{{}\today \hspace{.2in}\xxivtime}
\title{  % remember to have Vistula University!!
  Geographies of Wellbeing in Poland\\ (Urban-Rural, Regions, and Cities) 
}
\author{
% Adam Okulicz-Kozaryn\thanks{EMAIL: adam.okulicz.kozaryn@gmail.com
%   \hfill I thank XXX.  All mistakes are mine.} \\
% {\small Rutgers - Camden  % and Vistula University
% }
}

\begin{document}

%%\setpagewiselinenumbers
%\modulolinenumbers[1]
%\linenumbers

\bibliographystyle{/home/aok/papers/root/tex/ecta}
\maketitle
\vspace{-.4in}
\begin{center}

\end{center}


\begin{abstract}
\noindent
\end{abstract}
\vspace{.15in} 
\noindent{\sc %XXX TODO add to ebib as keyword PAPER-CODE-NAME and tag with ebib keywords 
}
\vspace{.25in} 

\vspace{-1in} 

\begin{spacing}{1.4} %TODO MAYBE before submission can make it like 2.0
\rowcolors{1}{white}{gray90}


{\huge[proposal after this]}

\section{changelog}
\subsection{oct17 met with piotr and incorporating his comments} {\large \color{red}srch 'TODO' in this sec}\\


\begin{spacing}{.7}
big/broad items
\begin{itemize}
\item  The document style requires improvement into more academic.--yes made it
  more academic, official and uptight TODO still rewrite text more, make uptight/acdemic
 \item added ad the end WPs/Tasks with KMs
 \item a
\end{itemize}
\end{spacing}

comments/notes responses
\begin{scriptsize}
\begin{spacing}{.89} 
\begin{verbatim}
Geographies of Wellbeing in Poland (Urban-Rural, Regions, and Cities)--bad title
geog, not soc
>>> rephrased away from geog to soc as i am going to soc and emphasize better
urb rur; i guess leave geog tit, but in writeup highlight sociology!!! it's
urban and rural sociology added wirth simmel toennies park and fischer; added
that the field is interdisciplinary and cited sociologists :) notably veenhoven
and delhey and bartram

emphasize better urb rur as important and overlooked, and definitely talk about
polska village wies! 
>>>added like there is little attention in the us and it has been left
behind add lang from my other cityAgency paper, we hypothetise it may be similar in pol

just using secondary data, not collecting my own data--definitely have to
highlight novelty and contribution!!!
>>>DONE highlighted whats  innovative at the beginning of  hypotheses sec 

show exactly how weak is current res and what is new here; cant be incremental,
need to be innovative; have more on problem statement--why is it important!!!
refer to currently going discussion and lit! using secondary data so explicitly
highligh how there is big ass contribution using these data!
>>>TODO more on problem statement--why is it important!!!

The research gap is not clearly presented.
Data is more described than the conceptual framework. Perhaps a new
conceptualization of urban rural well-being could prove a conceptual novelty.
>>>DONE again added to "hypotheses" section

The specific case of Poland with a context of the Second World War losses and a
particular development after the Second World War during the communist time can
be an interesting reference if the case is taken in a comparative perspective.
>>>WW2/communism another paper!; but here we go back to the 90s so post
communist development yes; and will compare to other eur c since 90s and their
development: added new hypotheis:  In a comparative perspective to other
European countries, the urban-rural SWB gradient is steeper in richer and faster 
growing economically countries. [EVS, WVS] 

Examining spatial and temporal dimensions of well-being dynamics seems
interesting. What are drivers of these dynamics--intermal or external (The EU and
the global economic processes, migration). The data is scarce.
>>>as in previous yes can do since 90s and v other eur c 

Poland experienced a dramatic increase of economic development measured by
GDP after 1990, but how it translates to well-being ins not clear.
>>>rt again can do 90s till now urb rur in pol v other eur c and gdp gro--and
already just descriptive just for poland is great--there was that easterlin
paper on happiness in economically rapidly growing china--do same for poland here
and add urb-rur :)

I would suggest to include the European comparative perspective as my colleagues
(P.Jabkowski, P.Cichocki) are very well trained and internationally recognised in
harmonisation of surveys data (https://doi.org/10.1093/poq/nfag064). It would be a
good justification of your location in AMU, Faculty of sociology.
>>>yes! again that H on gdp gro in eur c 
>>>TODO make sure in application form under motivation why AMU

The definition of well-being requires examination.
>>>i added in intro: "Both measures are defined and described in data
section."--there are pretty elaborate descriptions there--do we need more? what
exactly?

Suggested structure:
a) Background, problem description (substantive, theoretical/conceptual,
methodological) and objectives.
b) State of the art [=lit?]
c) Hypotheses, expected progress beyond the state the art.
d) Methods and data
e) Research plan (WPs)
f) Risks
>>>DONE! actually already kind of was done this way, just replaced sec names 
[i guess if it doesnt fly can redo more technical like mari did in her grants,
see her gmail "nie do publicznego wgladu"]

NAWA application btw has these:
3.6. Aktywnosci planowane w ramach Projektu: badania, prace rozwojowe, wdrozenia. Prosze opisac: 
stan badan; [=lit rev]
cele i zalozenia [=motivation, why this res]
pytania badawcze i hipotezy; 
proponowane narzedzia i metody badawcze; 
stopien innowacyjnosci;  = [lit contribution]
wspolprace z innymi instytucjami naukowymi lub przedsiebiorstwami, w tym
wspolprace miedzynarodowa. = [ebshoy and lina]



have some graphical conceptual model of analysis--piotr says super important
increases chances like 10% 
>>>see maris; maybe do fancy tikz--meh forget tiks, but added trancect with smiley faces

wnisoek pisac dla generalists
>>>TODO

background and objectives--czy temat jest warty podjecia!!
>>>TODO

US perspective: jest taki kraj polska co ciekawy; ale zmienic perspektyrw!!
>>>ok cut back a bit on that
porownac polske to innych w eruipe i do ameryki poludniowej! 
>>>ok done added on lina martinez and lat am 

and say while i do have us perspective being here 22 years, this is a strength,
i can compare us urb-rur to polish in this study; and again i am the leader of
the urb rur swb field, again that esseay in PIA sth from ru statements if not
double blind see what katrzyna says
>>>TODO


\end{verbatim}
\end{spacing}
\end{scriptsize}



specific items from comments in pdf
\begin{scriptsize}
\begin{spacing}{.29} 
\begin{verbatim}
style
>>>made it more uptight/academic

not convincing
>>>added references to sen and seligman

An US view. Simplistic.
>>>ha! yes we are simplistic here :) i cut back

Can be interesting but needs to be engaged in the current literature. Why
this gradient is relevant in EU, Poland, WOrld - needs justification.
>>>yes! done added soc lit; and elaborated on the gradient; and added diagram

It is very descriptive. WHat is the hypohteiss here: is Poland follow a
wider pattern with other Central Europe countries (perhaps a global
pattern?)
>>>no just that growth or shrinkage-->happiness--explained better
and more explicitly 

Good, but very descriptive. No breakthrough. A strong need for an
innovative conceptual framework.
>>>DONE added to hypotheses sec

Looks as territorilly embedded. Dat harmonization seems problematic. It
needs to be expained.
>>>added re harmonization

details needed.
>>>added new sec ``Work Plan: Key Milestones'' and refered the reader to there

Perhaps - discussion with the Catch-up lietarture of the Central EUrope? SOmehow
it is a default underlying concept.
POland compared to Other trasnsition coutries: Latin America? (See
https://books.google.com.fj/books?id=OBvQVhkIvdIC&printsec=frontcover
>>>added new hypothesis to compare with other eur countries; and added new hyopothesis on
social groups: age and education--the educated and younger ones may fare better
in cities v rural; and added a little bit on latin america/colombia too

Perhaps engagement with the energy transformation in Europe which requires
a massive change of consumption, would be important here as the large European
policy altering lifestyles and consequently well-being.
>>>yes!! added and citing my own paper on this!

It requires explanation why citing so old fashioned/obsolete
literature.
>>>ok dropped Fromm
\end{verbatim}
\end{spacing}
\end{scriptsize}


\begin{scriptsize}
\begin{spacing}{.19} 
\begin{verbatim}

\end{verbatim}
\end{spacing}
\end{scriptsize}

\section{Background, Problem Description, and Objectives}


The proposed study's outcome of interest is multidimensional human
subjective wellbeing (SWB) or flourishing. Two key dimensions are happiness or life
satisfaction\footnote{We use the terms interchangeably as is the usual practice
  in the social indicators literature, but specifically we mean global cognitive
 evaluation of life as whole rather than momentary and fleeing affective
 happiness.} and freedom (agency/autonomy).\footnote{Specifically we mean
 so called \texttt{freedom to}, not \texttt{freedom from} \citep{aok_free_from_to}.
Poles are \texttt{free from}: live in a free country, have no coercion,
 and are free from restrictions and impediments. But are they \texttt{free to}? 
 Are Poles actually able to choose, control and direct one's own life
 \citep{ee_ls15}?
% Not so much \citep{ee_ls15}.
}
 The measures are defined and described in data section. Wellbeing is
 an increasingly recognized and relevant metric in academic research, world
 development, policy and practice.   

The proposed project fits under the UN Sustainable Development Goals (SDGs): \#3  "Ensure
healthy lives and promote well-being for all at all ages." % and \#11 "Make cities
% and human settlements inclusive, safe, resilient and sustainable." %not really
% this one!
Happiness is now recognized as a development goal in
addition to GDP--the UN has passed resolution A/RES/65/309: ``Happiness: towards a holistic approach to
development'' stating that ``the pursuit of happiness is a fundamental human
goal'' and ``recognizing  that  the  gross  domestic  product  indicator [GDP]
by  nature  was  not  designed to and does not adequately reflect the happiness
and well-being of people.''\footnote{See
  \url{https://docs.un.org/en/A/RES/65/309}. Also see report promoting
  wellbeing as a human development measure by Nobel prize laureates Sen and Stiglitz \citep{stiglitz09al}.}

Freedom (agency/autonomy) is a human right as stated by the  Universal
Declaration of Human Rights (art. 26, 29): ``freedom for the realization of one's human potential.'' 
Freedom  is an end-in-itself%\citep{aok_free_from_to}
, as recognized by 
% \citet{fromm13,fromm12,fromm92,fromm64,fromm44,fromm94}
 Amartya   \citet{sen00}.

% \textbf{do check if this needs to be blind, then need to rewrite}
% I have coined the term ``urban-rural happiness gradient'' \citep{aok11a}
% % and the new line of literature has started.
% I am the leader of this field.\footnote{For instance see the discussion in framing essay IV in
%   \citet{johansen22}.}
There is a new geography of wellbeing field \citep[e.g.,][]{ballas2013geography}. 
 The field is interdisciplinary with key contributions from sociology such as
 Ruut Veenhoven
 \citep{brule2015geography,veenhoven1994satisfying,veenhoven2000freedom},
 Jan Delhey
 \citep{delhey2005predicting,delhey2003trusts,delhey2016happier,delhey2004life,delhey2002quality,alber2005quality}
 and David Bartram
 \citep{bartram2013happiness,bartram2011economic,hendriks2019bringing,international2013world}. 
The field has mostly international focus across
countries, with only few recent studies focusing on specific
countries (see next section).
 % TODO if katarzyna says ok then say i am teh leader of the field PIA essay etc
 % cp from ru statments


While earlier there was a need for focus on economic growth in Poland, now
 there is a need for focus on wellbeing. 
 In 2026 Poland has emerged as a one trillion dollar
 economy\footnote{\url{https://www.gov.pl/web/primeminister/poland-joins-the-trillionaires-club-a-historic-entry-into-the-worlds-top-20-economies}}% ,
 % and has converged closer to Western Europe
.  Polish PCGDP
PPP\footnote{Per Capita Gross Domestic Product Purchasing Power Parity adjusted.} 
is very close to that of Japan, and Polish minimum wage of 31 PLN or 8.15 USD
(as of 2026) is higher than the US minimum wage at 7.25 USD.  
%https://en.wikipedia.org/wiki/List_of_countries_by_GDP_(PPP)_per_capita
 %
% %It was not like that before, however.  
% I remember about 10 years ago I gave a talk on happiness in Warsaw, 
%%vistula, maybe little more recently like 2019
% and students didn't seem to see a need to care about happiness. Traditional
%topics like economic growth and personal income seemed more important. 
 %
% Wellbeing is increasingly relevant as Poland has been transitioning from
% developing country
% into a developed country so quality of life and
 Wellbeing is increasingly relevant over economic/survival basics such as income and
unemployment in the past. 
 When Poland was a developing country in the 90s, the focus rightly so was on survival
  and economic growth.  As Poland has
experienced economic miracle over past several years,\footnote{For instance Warszawski
Index Gieldowy (WIG) has increased about 120\% over past 5 years.} now it is
time to focus on wellbeing. 
 %
For visualization of the changing values from survival to wellbeing along with
the economic growth see figure \ref{ingle}.  

\begin{figure}[H]
 \includegraphics[height=2.1in]{/home/aok/papers/root/old/2025/swbColCitReg/tex/ingle.pdf}\centering
\caption{\label{ingle} From survival to wellbeing as economic growth increases
  \citep{inglehart97}. As Polish GNP/capita has increased
sharply over the recent years, the scholarly (and public) focus can and should shift from
survival to wellbeing.}
\end{figure}

Leading scholars advocate such change in focus from survival to flourishing. For
instance, Seligman \citep{seligman04,diener04s} led change in psychology that
traditionally used to focus on the negative side of human condition such as
 loneliness, anxiety, and depression--where psychologists have been trying to
 bring people from failure to normal. He pointed out that we should now focus on
 how to go from normal to flourishing. Amartya Sen is another example. Being an
 economist, he pointed out the limits of economic measures such as income and
 production, and proposed that the time has come to focus on wellbeing \citep{stiglitz09al,sen00}.
 %
 Poland still lacks a dedicated wellbeing study across its geographies and social
 groups. The proposed study aims to fill this gap.

\section{State of The Art: Urban-Rural, Regions, and Cities}

%!!!!!!!!!!!!!!!!!!!!!!!!!!!!!!!!!!!!!
%$>>>$guess cp here lit from my papers esp recent ones:\\
%(make sure i caover everything i have in hypotheses!!!)
%guess moree qol livability stuff and shrinking gorwing\\



%https://en.wikipedia.org/wiki/List_of_European_countries_by_area
%https://en.wikipedia.org/wiki/List_of_European_countries_by_population
% Poland is a medium-large country, 9th both in terms of area and population in Europe,\footnote{
% See \url{unstats.un.org} and \url{https://population.un.org}.}
  %
Poland is a diverse country in terms of its urban-rural, regional, and city geographies. 
%  And yet we know little about geographies of wellbeing in Poland and lack a
% dedicated comprehensive study. Hence, I propose this study to fill the gap.
 %
Urban-rural is not a binary, but a gradient--we will use several steps to
capture this gradient. Most scholars just use a binary, but wellbeing
raises nonlinearly from urban to rural intesyfying at population size of below several
hunderd thousand \citep{aok-ls_fisher16}. 

Urban-rural gradient matters a great deal, more than
most scholars realize. Indeed, it is difficult to overestimate its importance--the key
point is that urbanicity is not only a built environment, but a way of life
\citep{wirth38}. %--details in the data section.
 %
 %
Classic sociologists were aware of the profound impact or urbanism, but modern
empirical studies in this area do not properly ground their work in theory. 
The academic foundation of urban-rural wellbeing is sociological. Not only as 
urbanism as a way of life \citep{wirth38}, but also based on the greatest classic
writings in sociology: metropolis and mental life 
\citep{simmel03}, and Gemeinschaft and Gesellschaft \citep{tonnies57}.
 While there is a more recent empirical tradition in American urban sociology in study of urban wellbeing
% life notably by Park \citep{park15,park28,park84}, and more recently
 notably by Claude Fischer, this line of research has largerly been discontinued
 after 2000 \citep{fischer99,fischer95,fischer82,fischer81,fischer76,fischer75,fischer73,fischer72}. %fischer10,
% The literature has pointed to the following dowsides of urbanism: cognitive overloaded
% and overstimulation, vice, crime, and conspicuous consumption.
% urban  etiquette, sophistication, manners and finish  (Park, Veblen) bettencourt
 The proposed study aims to use the classic sociological theory with recent data
 to conduct a comprehensive study of the urban-rural happiness gradient in
 Poland. 
 
%Several urban-rural happiness gradient studies are worth mentioning. 
 \citet{weckroth21} is a closely related study to our proposed study--they analyze SWB by urban-rural
using ESS in Europe--we will also use ESS among other datasets, 
 but focus on Poland and add more recent waves and more wellbeing indicators. 
 %
\citet{knies26} is another close ESS study\footnote{And there are 2 more ESS
  studies: \citet{ivanovic2022subjective} is an urban-rural study, but again
only 2 waves and no focus on Poland. \citet{vaznoniene2019wellbeing} is
interesting ESS study in a sense that it does focus on one country and discusses
country specific intricacies like us, but it is Lithuania. \citet{aok21} is a major recent study using WVS, but as most research it focuses on cross-country
 patterns, not Poland.} --it uses multiple indicators, but only  wave 6 and like other
research does not focus on Poland. %We will use all 12 waves.
 %
There are no studies on urban-rural freedom gradient, not just in Poland but
anywhere--ours is the first.\footnote{\citet{yamagishi12} studies urban-rural
  freedom but in very different way from us--they show that urbanites are more
  likely to chose extravagant pens than rural folks. We will study freedom for
  urban-rural and regions (we do not have 
  freedom variable for cities).}

From another perspective, there is literature on ``left behind ruralities.''
Rurality in the developed world is increasingly a periphery, loosing power and voice v
 urban centrality, being abndoned and disrespected--there is rural disconent,
 resentment, and shame. See superb accounts in 
 \citet{hansonCityJournalautumn15,fullerNYT17monD}.
 Rural areas are 'left-behind' \citep{ulrich2021people}, 'places that
don't matter' \citep{rodriguez2018revenge}%cited like 4k!!
---they are second-tier areas with limited or deteriorating opportunities.
% Ruralities feel that policy is done to them, not with or for them. 
There may be similar pattern in Poland and it may work to decrease rural SWB
that otherwhise should be higher than urban.
 %
Next we turn to regions.
 
 Regions are an useful level of analysis as there are
 great differences within countries and across regions, and Poland is no
exception.  There are life satisfaction studies across regions, but no study focuses on Poland \citep[e.g.,][]{aok10_geo,ballas23,ziogas23,ziogas20}.
Similarly, quality of life and wellbeing studies across regions and cities focus
on cross-country patterns, not Poland, and tend to use older datasets \citep{aok-swbLivability18,aok13liavbility}.
Poland has 16 voivodeships, so called NUTS2 level
regions.\footnote{Voivodeships$=$NUTS2, except  Mazowieckie (the Warsaw
  region)--1 single voivodeship, but 2 NUTS2: Warsaw Metropolitan and Mazovian
  Region; so total of 17 NUTS2 regions.} %(in European statistical terms).
 At NUTS2 level we will investigate geographic clustering and the relationship
 between objective development measures such as income and subjective measures (SWB). 
 %
There are also opposite population dynamics at NUTS2 level, for
instance,  Mazovia (where Warsaw is located) has been growing, but most  
other regions have been stable or shrinking. We will investigate the effect of
population change on wellbeing--we know that population dynamics are highly
correlated with SWB \citep{growShrink25}, but again, we do not know the
relationship in Poland. The hypothesis is that population growth predicts
greater wellbeing, and population shrinkage predict lower wellbeing. Population
growth signifies bright outlook, and decline signifies difficult times coming. 
The theory that may explain the mechanism between population change and SWB is  the so-called ``tunnel effect.'' Humans tend to think that
whatever is happening around, whether things are getting better or worse,  will
eventually happen to them as well.\footnote{
 \begin{quote}
  Suppose that I drive through a two-lane tunnel, both lanes going in the same direction, and run
  into a serious traffic jam. No car moves in either lane as far as I can see (which is not very
  far). I am in the left lane and feel dejected. After a while the cars in the right lane begin to
  move. Naturally my spirits lift considerably, for I know the jam has been broken and that my
  lane's turn to move will surely come at any moment now. Even though I still sit still, I feel much
  better off than before because of the expectation that I shall soon
  be on the move \citep[Hirschman, quoted in][p. 88]{ravallion00l}.
\end{quote}% \vspace{-.12in}
}
%https://en.wikipedia.org/wiki/List_of_cities_and_towns_in_Poland

Finally, we have data representative of 
Bialystok, Gdansk, Cracow, and Warsaw with domain satisfactions (life, place,
personal, and financial), and various quality of life indicators regarding
amenities (e.g., public goods and services such as parks)  and
disamenities (e.g., community problems such as crime). We will investigate 
these across the cities.    
 %
 % For instance, Warsaw has been clearly growing, while other cities such as
 % Poznan have been mostly flat, and yet others such as Lodz and Szczecin have
 % been shrinking.
 %
 Specifically, large Warsaw (and perhaps Cracow) is likely to differ from small
 Bialystok. 
 Warsaw stands out as not only the largest city, but also possibly a primate
 city, disproportionately larger and more influential than any other city.\footnote{At
$>3m$ Warsaw metro is larger than Katowice metro at about 2.5m and third
Karkow metro at 1.5m.
 %
While Katowice metro is not half or less of Warsaw metro as
 ``the law of the primate city''
 would indicate \citep[see discussion in][]{aok12zipf}, it is only due to the
 location of Katowice in regional industrial cluster (Gornoslaski Okreg
 Przemyslowy (GOP)), Katowice city itself is much smaller than  Warsaw at
 0.28m v 1.9m. 
 %
Maybe Warsaw is not such clear primate as Cairo  \citep{ebshoyEgypt24}, but it
could be similar in its primate status to Bogota \citep{unhappyBog26}.
}
 
Last but not least, there are several related Polish studies in smaller local
journals that are worth mentioning as they are mostly about Poland
specifically--again, the international literature does not focus on Poland.
 %
 We start with a hidden gem: \citet{cichocki2009changes}--this is exactly the
 kind of study we want to conduct--just using more and newer data, controling
 for relevant predictors of SWB, and situating the study in the international
 literature. 
 %The closest first.
 \citet{sompolska2020quality} is quite close, but uses different data: Polish
Social Diagnosis dataset and does not control for SWB predictors. 
 %
\citet{vaznoniene2021territorial} uses ESS and studies urban-rural SWB like us, but
 focuses on Lithuania v Poland, and uses waves 4 and 9 only and 3 urbanicity
 categories only, and does not control for SWB predictors. \citet{wojewodzka2023rural}
 is the same, just uses waves 2 and 8. 
  %
\citet{tobiasz2017urban} does the analyses across the regions and urban-rural,
 %and while social connection and loneliness are closely related to SWB,
 but their sample is elderly and the study is about health only.
 \citet{stangierska2022urban} studies quality of life across green v not green
 urban areas, but only across Warsaw neighborhoods, and the sample size is small at 381 persons.
 \citet{sompolska2021happier} uses Polish Social Diagnosis and studies urban-rural 
 SWB--but separately: predictors of SWB in urban separately from predictors of
 SWB in rural. %--it is different from our study.
 Overall, while there are studies focusing on Poland, like our proposed study, they have
 important limitations as summarized above.\footnote{ 
 And there are several studies that differ from the
 above. \citet{churski2022geographical} has very fine 
 geographical resolution of 380 units across Poland, but they do not study SWB,
  only the related objective indicators such as jobs. 
 \citet{weziak2025understanding} is one of the very many of the studies
 co-authored with VanderWeele, this time in Poland, but there is no geography in
 the study.  
 %
 \citet{wkeziak16} uses the same flash barometer 366 that we will use, but does
 not focus on Poland and instead on all European cities. 
}
 Hence,  we know little about wellbeing geographies in Poland--we
 turn to our study.%, yet un-answered, hypotheses.  

\section{Innovative Conceptual Framework And Hypotheses--Expected Progress Beyond the State the Art}

% The research question: 
% This research is about geography of wellbeing in Poland (urban-rural gradient, regions, and cities).

% The hypotheses listed below are a primary set, we will be refining them and
% adding other hypotheses at the literature review stage.

\noindent %Hypotheses
%(SWB$=$(multidimensional) Subjective Wellbeing)):\\

%The research gap is great, and the proposed study clearly fills the gap.
It will be the first such study in Poland.
 %
The
classic scholars of 19th and 20th century put forward research theories and
theses that could not be properly tested at the time due to lack of data and computing.
We ground our research better in classic theory, which allows us to pose research
questions overlooked by the contemporary research. 
 %
% opis zjawiska co ja badam, jaka jest innowacja jak
%nowo inaczej podchodzimy; wusze cos wymyslic; od razu pierwsze jest ze better
%measurement of urb-rur as gradient; more measures of swb like from ess etc
 %
The innovative conceptal framework includes elaborated urban-rural gradient
measurement.
 We will use 5-8 bins depending on the survey %ess5 wvs8
to measure urban-rural gradient, a critical 
improvement over usual binary or ternary urban-rural measurement in the literature.
 Urban-rural wellebing gradient is visualized in figure \ref{urbRur}.

 \begin{figure}[H] %like wykretowicz dissertation! conflict potrzeb cywilizacyjnych!
%  \includegraphics[height=5.5in]{Humanwildlife_conflict_labeled.png}\centering
    %From left to right: urban-rural. The middle panel from left  (human icon) to righ (deer icon) symbolizes
% Top panel is pre-sprawl and bottom is post-sprawl--we
    % will capture sprawl with suburbanization variable at person level; %ess 
    % and indirectly at regional and city levels.
    % flash barometers! for regions and cities 
    %Source: \url{https://commons.wikimedia.org/wiki/File:Humanwildlife_conflict_labeled.png}
    %\includegraphics[width=5in]{aok.png}\centering
    \includegraphics[width=5in]{transect.png}\centering
    \caption{\label{urbRur} From right to left:
      Urban-rural gradient visualization with transect diagrams symbolizing
      urban-rural environments and corresponding frownie to
      smiley emojis symbolizing wellbeing gradient.
    From urban to rural there is increasing nature, biodiversity and
    hypothetiszed wellbeing. Source: Authors. The transect
    diagram for urbanicity is by Duany Plater-Zyberk \& Company
    \url{https://transect.org/rural_img.html}}
  
 \end{figure}

 Another innovation is a more extensive measurement of wellbeing.
We will use a battery of survey items to capture  multidimensional wellbeing capturing not only life satisfaction
as most literature does, but also functionings, agency, and vitality.  
 %
 Combining multistep urban-rural wellbeing gradient with multidimensional
 measurement of wellbeing and focus on Poland, our study is the first of its
 kind--nobody to date has used such a research design.
 %
 Finally, our innovative conceptal framework includes using large multiple
 datasets at multiple levels: urban-rural (and over time since the 90s),
 voivodeship (NUTS2), and city (see next data section).
 %
 The proposed research will test the following hypotheses to produce new
knowledge. The new knowledge is not only desirable on its own (academic value),
but it also has policy and practical implications--see section ``Significance,
Policy, and Practice.''


\textbf{Q ??:} too many H? cut some? which ones? on the other hand this is for 4-5
  papers over 4yrs so maybe ok? \\

\noindent $H_1$ Cities have lower multidimensional Subjective Wellbeing (SWB) than rural areas, with the largest
drop occurring for the largest places. [WVS, ESS, FB539]\footnote{WVS$=$World Values
  Survey; ESS$=$European Social Survey; FB539$=$Flash Barometer \#539--these and other datasets are described
  in the next section.}\\ %ess
\indent$H_1a$ Certain social groups such as working-age adults and higher social classes
have higher SWB in cities v rural areas. [WVS, ESS]\\
\indent$H_1b$ Over time, from 1989 to 2012, the SWB gap between cities and rural
places has increased. [WVS]\\ 
\indent$H_1c$ In a comparative perspective to other European countries, the urban-rural
SWB gradient is steeper in richer and faster growing economically countries. [EVS, WVS]\\
$H_2$ More remote, less populated, and lower density regions have higher SWB than
more central, more populated and higher density regions. [ESS, FB539]\\ %flash b: just maps no regressions too few units
\indent$H_2a$ Regions that experienced population growth will have higher SWB than flat
or shrinking regions. [ESS, FB539]\\ 
\indent$H_2b$ Across regions correlation between objective quality of life measures such as income with
SWB measures will be moderate only at $<0.5$  thus revealing inadequacy of
  traditional progress/development measures such as
income. [ESS, FB539]\\
\indent$H_2c$ Urban-rural SWB gradient across regions will be relatively stable,
i.e., larger places have higher SWB than smaller places. [ESS, FB539]\\ 
$H_3$ The relationship between quality of life (local amenities and disamenities) and domain satisfactions (life, place, personal, and financial) will differ across                                  
 Bialystok, Gdansk, Cracow, and Warsaw. [FB366,419] %\citep{aok-swbLivability18}

 
 
\section{Data, Measurement, and Methods}
%maybe that diagnoza spoleczna--or may turn out its crap, maybe someone already said so

%\subsection{Multidimensional Wellbeing}
The goal is to do a comprehensive study of Polish geographies of wellbeing, hence
we use multiple datasets, each having its own strengths: World Values Survey (WVS), European Social Survey
(ESS), and Flash Barometers (FB) for regions (FB539) and cities (FB366,419).

WVS has a usual SWB item:
 ``All things considered, how satisfied are you with your life as a whole these days?''
 on scale 1-10. 
Such SWB measures tend to be reasonably valid and reliable
%  \citep{myers00,ditella06m,layard05,bray06g,sandvik93ds,clark08al,diener09}.
\citep[][ch. 5]{diener09}. 
% 
WVS allows us to examine the geographies of wellbeing over time, which is important
  as Poland has evolved greatly, and WVS data goes back to 1989. We have 4k observations, about
  1k per wave.\footnote{We use waves 1989, 1997, 2005, and 2012.
 For the 1997, 2005, and 2012 waves we have 8
rural-urban categories -2k, 2-5k, 5-10k, 10-20k, 20-50k, 50-100k, 100-500k, and
500k-. In 1989 we only have information about 49 voivodeships, but
 urbanicity can be inferred from this information indirectly,
 and proximity to a major city directly.} 
 Also, an important advantage of WVS is
 that urbanicity is recorded by an interviewer.\footnote{It is more objective than subjective respondent self-reports. For
   instance, respondents may not know well whether they live in a 20k or 30k place.}

Freedom is measured in WVS,\footnote{"Some people feel they have completely free choice and control over their lives, while other
      people feel that what they do has no real effect on what happens to them. Please use this scale
      where 1 means 'none at all' and 10 means 'a great deal' to indicate how much freedom of choice
      and control you feel you have over the way your life turns out." 
    Validity and reliability of the measure seem fine and the variable has been used successfully
 in the literature \citep[e.g.,][]{verme09,inglehart08,veenhoven00_b,aok_free_from_to}. }
%if anything it could be interpreted slightly differently in different
%languages/cultures but that is not a problem here for Poland only.}
 and in some ESS waves 
 %which has several waves %3,6,12
 containing a whole battery of wellbeing items measuring functionings and vitality. 
 %
Preliminary analyses using ESS 1-6 cumulative file already had over 10k
observations for Poland.\footnote{The data will be updated upto wave 12 that will be
available in 2027.}
  %
 We will use ESS items such as those in wave 6 on functioning ``free to decide how to live my life,''
  %``enthusiastic about what you are doing,''
 and vitality ``had lot of energy.''\footnote{ESS urbanicity is measured with: a big city, suburbs or outskirts
of big city, town or small city, country village, farm or home in
countryside. Per wellbeing see overview at 
  \url{https://www.europeansocialsurvey.org/findings/europeans-wellbeing/measuring-wellbeing}. Wave 
6 wellbeing items:
\url{https://europeansocialsurvey.org/sites/default/files/2026-01/ESS12-wellbeing-module-design-template.pdf}
 see Table 1: Items from the wellbeing module grouped by the dimension of
 wellbeing they relate to. We also will use new wave 12 that will be available in 2027 and will
 contain a wellbeing module. See \url{https://europeansocialsurvey.org/sites/default/files/2026-01/ESS12-wellbeing-module-design-template.pdf}.}
 

FB539 (2024) has 5.2k observations and is representative of
 17 NUTS2 regions, and has a usual SWB
 question similar to WVS.\footnote{We may also use
  the earlier flash barometers: \#497 (2021), \#472 (2018), \#427 (2015),  \#356 (2012)
  depending on how strong and robust the latest wave results are.} 
% \begin{figure}[H]
%  \includegraphics[height=2in]{nuts2.png}\centering
% \caption{\label{nuts2} 17 regions (voivodeships) so called NUTS2 regions}
% \end{figure}
 %
FB366,419 (2012, 2015) are representative of 4 Polish cities:
 Bialystok, Gdansk, Cracow, and Warsaw. The sample size is about 2000
 observations per wave, and again we have a usual SWB question like in WVS.
%later can also try Flash Eurobarometer 156 (Urban Audit Perception Survey 2004)

Data harmonization is not needed in WVS--one advantage of the survey is
consistent measurement over time. We will use the data
file  %so called ``time-series''
that has all the waves in it with consistent coding of the variables.  FB539 is
a cross section. FB366,419 are 2 repeated cross sections (2012, 2015)--we will
check coding of the survey items and if they differ and if we make over time
comparisons using these data, we will harmonize. The same pertains to ESS.  

\subsection{Models}


We will mostly use ordinary least squares (OLS).\footnote{Although OLS assumes cardinality of the
outcome variable, and SWB measures are technically ordinal, OLS is an appropriate estimation method.
\citet{carbonell04} has shown that OLS results are substantially the same as
those from discrete models, and  OLS has become the default method in happiness
research \citep{blanchflower11}. Theoretically, while there is still debate
about the cardinality of SWB, there are strong arguments to treat it as a
cardinal variable \citep{ng96,ng97,ng11}.} For robustness/sensitivity
checks for the key models we will use multinomial logit (as the ordinal logit
tends to fail proportional odds assumption). As another robustness/sensitivity
check we are considering machine learning models such as gradient boosting
models \citep{oparina25}, specifically XGBoost
 (\url{https://xgboost.readthedocs.io}). 

%\section{controls}
We will use a usual set of controls such as those in the literature: 
for life satisfaction \citep[e.g.,][]{aokcities,aokrel}; and for freedom
\citep[e.g.,][]{pitlik2016free,nikolaev2016give}. 
 %
 A key robustness/sensitivity check will be specification curve analysis (SCA)
\citep{simonsohn20}. And %, related,
 we will pay extra attention to the choice 
 of control variables and avoid overcontrol bias.\footnote{
%\begin{quote}
``What we really need, to identify/estimate X$->$Y, is to control for other
determinants of Y that are also antecedents of X (so, W$->$X). If we include
controls where the relationship goes the other way (X$->$W), we will exacerbate
bias in our estimate of X$->$Y. (The relevant term for that situation is
'overcontrol bias')'' \citep[][p. 5]{bartram24}.
%\end{quote}
}

To examine geographic patterns among regions we will use choropleth maps,
Moran's I, and Local Indicators of Spatial Association (LISA) \citep{anselin1995local}.\footnote{ 
We may do some multilevel modeling (MLM) at NUTS2 level as in
\citet{weckroth21} but focusing on Poland. Yet MLM is problematic with only 17
U/As at Level 2--if we use it, results will be postponed to appendix. 
We may use Small Area Estimation \citep{moretti24}. 
}



\section{Significance, Policy, and Practice}

%guess third of points is: Potencjalne znaczenie Projektu dla rozwoju nauki

The proposed research will contribute to the literature and produce new
knowledge.   
 We will write at least 4 peer-reviewed articles\footnote{Out of which at least
   3 will be published in major international journals such as Social Indicators
 Research, Applied Research in Quality of Life, and Cities, with IF$>$3.} in this project, and through further research
mentoring and leadership additional ones in the long term (see section
``3.15. Didactic/mentoring activities'' in the application form). Regarding the
 activities during the duration of this project see in this document section ``Work Plan: Key
 Milestones.'' 
%See for details in the application form under ``Dzialania mentoringowe.'' 
%\textbf{see if thats where it fits}

The research also has policy and practical implications. First, as stated in the
introduction, the time has come for Poles to focus on wellbeing as opposed to
just survival. Poland has modernized infrastructure, higher incomes, and booming economy. The %past 
focus on survival and economy
 can now give a way to 
 %now shift towards
 wellbeing and flourishing.
  With the proposed study we aim to aid in this pursuit of happiness. 
 This is called evidence-based pursuit of happiness in the literature \citep{veenhoven12}.\footnote{There is a
  psychological nuance here--one should not focus fully on the pursuit of
  happiness and chase it or value it intensely
  \citep{mauss11,gruber11}, neither be concerned about it and judge it
  \citep{zerwas24}--it may backfire. Rather, one needs to know what is important
for wellbeing from the academic literature, like the proposed study, and include
this information in everyday decision-making. For instance, humans should
 buy experience such as travel and time with others, not things such
as conspicuous consumption items \citep{aok_ls_car15,veblen05a,veblen05b,thoreau95}.}
 The proposed study will find out which parts of Poland enjoy high wellbeing and for
 whom--which social groups, e.g., by age and social class, flourish where? For
 instance, it may be the young and the educated that do not suffer from urban SWB
 penalty \citep{aok14}.   
 Our findings should help Poles increase their wellbeing.

There are many avenues for specific policy and practice collaborations.
{One simple implementation is  to educate policy
 makers, public administrators, and citizenry about the wellbeing metrics and
 study's findings so that we can make more informed decisions. 
 %
 To that end we will share and popularize ourfindings through public lectures, 
  popular press, and social media. In the long term we are thinking about
  establishing a wellbeing research center in Poland, similar to Erasmus
  Happiness Economics Research Organisation (EHERO; \url{https://www.eur.nl/ehero}), where PI used to work. 
   
  
 SWB can and should be used for public policy and administration \citep{diener09}.
% provides a good discussion of why potential problems with happiness are not serious
% enough to make it unusable for public policy--see especially  ch. 6.  
% These potential problems include  genetic
% determination of happiness \citep{lykken96t},  adaptation
% \citep{brickman78cj},  and various comparisons \citep{michalos85}.}  
% 
 SWB can serve as a ``yardstick'' \citep{veenhoven88}--arguably the ultimate outcome of public
 policy and administration is SWB.\footnote{While the idea of using wellbeing data for policy and
    administration is new in Poland, it is not new in general
    \citep{stiglitz09al,diener09}.} "If you can't measure it, you can't improve
 it"--if governments don't measure wellbeing, they can't improve wellbeing.  
  I will reach out to Polish government to implement the study findings in
  practice. 
  I plan to start locally with Wielkopolska voivodeship,
  where I will be physically carrying out the study, and other regional
  governments may be involved as well. %as I will be carrying out the study.
  Likewise, some city
  governments may be involved as the study will use data representative of 4
  cities.


  It is relatively straightforward to implement the study's findings. For instance,
  we will find where and which social groups flourish and languish--such
  information can be then used to advance utilitarianism, produce the ``greatest happiness for the
  greatest number'' \citep{akay15,nussbaum12,kahneman97ws}.
  %
  Another avenue is to use the wellbeing metrics to evaluate policy. For instance, if
  there a major change in policy or investment in public goods, we can actually
  measure whether it produced an increase in wellbeing by using before-after
  comparative group research design (difference-in-difference).} We could use a
SWB yardstick to measure policy effectivness in terms of SWB. For instance,
we identify a problem, say income inequality--see figure \ref{fOL}. We can
 take action with progressive taxation to achieve redistribution that would
 produce outcome of interest, income inequality. And we usually stop there, but
 we should probe further--have we improved human lot, did we improve quality of
 life, has wellbeing increased? Usually\footnote{There are few
  exceptions to this rule, and SWB is not the only outcome
 of interest especially in the short run, indeed, we may often 
 need to sacrifice SWB in the short run in order to achieve it in
 the long run--for instance save, invest, postpone consumption, and work
 hard in order to increase SWB later.} the outcome of interest is only
worthwhile if it leads to the ultimate outcome of interst, SWB \citep{aok_lsPol16}.

% Define block styles
\tikzstyle{s} = [rectangle, draw, fill=black!10
     % text centered
    , rounded corners,  text width=2.8cm]  %node distance=2cm, minimum height=0em,  minimum width=2cm,
%\tikzstyle{s} = [draw, thin, fill=blue!20] 
\tikzstyle{e} = [rectangle, draw,  fill=black!10,   text width=2.8cm, , rounded corners  %, minimum height=2.5em
% , text width=5em
]
\tikzstyle{block} = [rectangle, draw, fill=black!20, 
    text width=10em, text centered, rounded corners, minimum height=4em]
\tikzstyle{b} = [rectangle, draw,  
    text width=6em, text centered, rounded corners, minimum height=4em]
\tikzstyle{line} = [draw, -latex']
\tikzstyle{cloud} = [draw, ellipse,fill=black!20, node distance = 5cm,
    minimum height=2em]
   
    \begin{figure}[H]
      \centering
      \begin{tikzpicture}[thick, node distance=3.7cm, text centered
        ,font=\small]
        % Place nodes
        \path[->] node[e,] (prob) {problem \\ (income inequality)};
        \path[->] node[s, right of=prob] (act) {activity \\
          (progressive taxes)} (prob) edge node {}(act);
        \path[->] node[s, right of=act] (sub) {action/subobjective \\
          (redistribution)} (act) edge node {}(sub); \path[->] node[s,
        right of=sub] (out) {outcome of interest \\(income
          equality)} (sub) edge node {}(out); \path[->] node[s, right
        of=out] (hap) {ultimate outcome \\ (\textbf{SWB})} (out)
        edge [dashed] node {}(hap);
      \end{tikzpicture}

    \caption{An outcome line. Usually we stop at {outcome of
        interest}. Yet the  {outcome of
      interest} is only of value (with few exceptions) if it results in SWB.}
\label{fOL}
\end{figure}

One avenue would be to engage with energy policy. While in general energy use
has been found to be unrelated to SWB earlier \citep{okulicz19}, this claim 
needs to be re-examined amidst current energy crisis and skyrocketing
prices. Massive reductions in consumption and altering of the lifestyles may
ultimately reduce wellbeing. 
 % 
  Finally, long term (after this project) it could be data
  collection, say implementing a survey similar to Flash Barometers locally; or
  using social media data to meeasure wellbeing \citep{iacus25b}. 
 
 In framing policy and practice discussion we will also use Latin American
 perspective. One of the team members is a senior wellbeing scholar from
 Latin America, a public policy researcher whose work integrates the analysis of
 social programs, impact evaluation, and subjective well-being metrics. Our
 expert has also practical experience--has served as a consultant for the World Bank and the Inter-American
  Development Bank in Latin America. Latin America is the happiest region in the
  World \citep{rojas15,rojas25}, and while there are similarities
  between Poland and Latin America, such as high levels of social connection,
   there are specific actionable lessons from the region that other countries
   can use \citep{colQolSwb25}. 
  
  
\section{Research Plan: Work Packages (WPs) and Tasks (Ts); and Work Plan: Key Milestones (KMs)} 
%Key Milestones Plan prac: kluczowe kamienie milowe dodac do wniosek 3.12

\begin{spacing}{.7}
note: in square brackets are enumerated team members responsible for each task;
order indicates workload from high to low
\begin{itemize}
\item all$=$all team members
\item aok$=$adam okulicz kozartn (principal investigator) 
\item lm$=$lina marinez
\item em$=$ebshoy mikhaeil
\item pm$=$piotr matczak ???see katarzyna email and if ok TODO add to tasks
\item mb$=$mariusz baranowski
\item pj$=$piotr jablkowski
\item pc$=$piotr cichocki
\item student$=$ AMU (Adam Mickiewicz University) graduate student; there are 3 of them
\end{itemize}

\begin{itemize}
\item WP0 monthts 0-48: management [aok] 
  \begin{itemize} 
  \item T0 months 0-48 management: project management and coordination  notably tracking the timeline
    and managing progress, communication, and risks [aok]
  \item T1 months 0-2 meet and greet: conduct a series of in-person and zoom
    meetings to plan and get going with the project [all] 
  \item T2 months 0-2 setup project: complete necessary paperwork, notably employ team members at the
    university [all]; note: IRB approval is not necessary as we are using
    secondary data only
  \item T3 months 0-48 meetings: in-person (whenever possible) or zoom (if necessary) weekly or at very
    least every two weeks meetings of team members currently involved in project
    tasks with PI [aok, team members currently working on tasks]
  \item T4 months 45-48 report and finalize: write final report and finalize the project [aok,lm,em]
  \end{itemize}
%\end{itemize}
\item \colorbox{black}{\textcolor{white}{KM0 month2: setup done}}
%\begin{itemize}
\item WP1 monthts 2-12: literature review
  \begin{itemize}
  \item T1 literature review: assign literature domains to 2 AMU students, and have students
    perform the review [2students]
  \item T2 literature review supervision: have more senior team members 
    supervise the students and
    hold periodic meetings at least every 2 weeks or as necessary [em,aok]
    % !!use AI to check if we missed any literature
  \item T3 finalize literature reviews: draft and revise literature reviews for
    subsequent 4-5 journal articles [2students,em,aok] 
  \end{itemize}
\item \colorbox{black}{\textcolor{white}{KM1 month12: literature review done}}
\item WP2 months 2-18 H1 papers: data management, analysis, and 2-3 papers drafting for hypothesis H1
  \begin{itemize}
  \item T1 data management: clean, organize, document, recode [em,1student,aok] 
  \item T2 exploratory data analysis: descriptive statistics and visualizations [em,1student,aok] 
  \item T3 main analyses: mostly regressions and robustness checks (see for
    details methods section) [em,aok,1student]
  \item T4 analysis review and revisions: review and revisions of the analyses by the senior team members [aok,pc,mb,pj]
  \item T5 paper drafting and revisions: drafting of the papers and revisions [aok,lm,em,2students] 
  \end{itemize}
\item \colorbox{black}{\textcolor{white}{KM2 month18: data analyzed and 2-3 papers drafted for H1}}
\item WP3 months 18-30 H2 paper: data management, analysis, and 1 paper drafting for
  hypothesis H2; Tasks are the same as in WP2
\item \colorbox{black}{\textcolor{white}{KM3 month30: data analyzed and 1 paper drafted for H2}}
\item WP4 months 30-38 H3 paper: data management, analysis, and 1 paper drafting for
  hypothesis H3; Tasks are the same as in WP2 
\item \colorbox{black}{\textcolor{white}{KM4 month38: data analyzed and 1 paper drafted for H3}}
\item WP5 months 18-48 dissemination: academic dissemination, real-world popularization, and policy engagement
  \begin{itemize}
  \item T1 articles:  $>=$ 3 submissions to journals with $IF>3$ such as
    \texttt{Applied Research in Quality of Life}, \texttt{Social Indicators
      Research}, or \texttt{Cities} [all]
  \item T2 conferences: 3 conference presentations such as  \texttt{International Society for
      Quality-of-Life Studies (ISQOLS)} or \texttt{European Regional Science
      Association (ERSA)} [aok,lm,em]
  \item T3 lectures:  $>=$4 public lectures and or seminars at AMU and or other Polish universities
  \item T4 policy engagement: engage with local public policy and administration starting in Poznan
    and Wielkopolska [aok,1student]
  \item T5 popularization: popularization via popular media (e.g., Glos Wielkopolski) and or social
    media (e.g., YouTube) [aok,lm,em,1student]
  \end{itemize}
\item \colorbox{black}{\textcolor{white}{KM5 month48: dissemination and popularization done}}
\end{itemize}
\end{spacing}
% draft grant applications for subsequent funding. We are considering the following: NCN Opus,
% EU Horizon, and ERC grants, but there may be other opportunities as well. submit grant applications 


\begin{ganttchart}[
  x unit=0.3cm,
  vgrid= {*{1}{gray, dotted}}, %, *1{black, dashed}
   bar label node/.append style={
     align=left,
     text width=width("WP5 dissemination")}
   ]{0}{48}
\gantttitle{Year 1}{12}
\gantttitle{Year 2}{12} 
\gantttitle{Year 3}{12}
\gantttitle{Year 4}{12}\\
\ganttbar{WP0 management}{0}{47} \\ %ending val need to substrat 1
\ganttbar{WP1 lit rev}{2}{11} \\
\ganttbar{WP2 H1 papers}{2}{17} \\
\ganttbar{WP3 H2 paper}{18}{29}\\
\ganttbar{WP4 H3 paper}{30}{37}\\
\ganttbar{WP5 dissemination}{18}{47}\\
\end{ganttchart}


\section{Risks}

see application form ``3.9. Risk analysis and security issues''


%  instead \ExecuteMetaData[../out/tex]{ginipov} do \emd{ginipov}

% \begin{figure}[H]
%  \includegraphics[height=3in]{../out/gov_res_trust.pdf}\centering\label{gov_res_trust}
% \caption{woo}
% \end{figure}


%TODO !!!! have input here aok_var_des



% %table centered on decimal points:)
% \begin{table}[H]\centering\footnotesize
% \caption{\label{freq_im_god} importance of God}
% \begin{tabular} {@{} lrrrr @{}}   \hline 
% Item& Number & Per cent   \\ \hline
% 1(not at all)&    9,285&  9\\
% 2&    3,555&        3\\
% 3&    3,937&        4\\
% 4&    2,888&        3\\
% 5&    7,519&        7\\
% 6&    5,175&        5\\
% 7&    6,050&        6\\
% 8&    8,067&        8\\
% 9&    8,463&        8\\
% 10&   52,385&       49\\
% Total&  107,324&      100\\ \hline
% \end{tabular}\end{table}


% % Define block styles
% \tikzstyle{block} = [rectangle, draw, fill=black!20, 
%     text width=10em, text centered, rounded corners, minimum height=4em]
% \tikzstyle{b} = [rectangle, draw,  
%     text width=6em, text centered, rounded corners, minimum height=4em]
% \tikzstyle{line} = [draw, -latex']
% \tikzstyle{cloud} = [draw, ellipse,fill=black!20, node distance = 5cm,
%     minimum height=2em]
    
% \begin{tikzpicture}[node distance = 2cm, auto]
%     % Place nodes
%     \node [block] (lib) {liberalism, egalitarianism, welfare};
%     \node [block, below of=lib] (con) {conservatism, competition, individualism};
%     \node [cloud, right of=con] (ls) {well-being};
%     \node [block, below of=ls] (cul) {genes, culture};
%     \node [b, left of =lib, node distance = 4cm] (c) {country-level};
%     \node [b, left of =con,  node distance = 4cm] (c) {person-level};
%     % Draw edges
%     \path [line] (lib) -- (ls);
%     \path [line] (con) -- (ls);
%     \path [line,dashed] (cul) -- (ls);
% \end{tikzpicture}


%PUT THIS NOTE, polish and put to /root/author_what_data --ALWAYS
%stick here stuff as i run it!!! maybe comment out later...


%!!!!!!!!!TODO: see cityAgency for cool setup :)))

% \section*{\Huge ONLINE APPENDIX}
% \textbf{[note: this section will NOT be a part of the final version of
%   the manuscript, but will be available online instead]} %hence everything below
%                                 %is organized byu section, not subsection
% !!!
% have most of the stuff outputted to online appendix:)--start with that and then
% select stuff to paper--have brief narrative describng patterns in online app too
% !!!

% \section*{Variables' definitions, coding, and distributions}
% \label{app_var_des}


% %\input{/tmp/a.tex} %aok_var_des

% % \begin{spacing}{.9}
% %   \begin{table}[H]\centering \caption{Summary statistics.} \label{sumSta} \begin{scriptsize} \begin{tabular}{p{1.8in}p{.5in}p{.5in}p{.5in}p{.5in}p{.5in}p{.5in}p{.5in}p{.5in}p{.5in}p{.5
% %             in}p{.5in}p{.5 in}}\hline
% %         \input{/tmp/aha2.tex}
% %          \end{tabular}\end{scriptsize}\end{table}
% % \end{spacing}

% % \begin{spacing}{.9}
% %   \begin{table}[H]\centering \caption{Correlation matrix.} \label{sumSta} \begin{scriptsize} \begin{tabular}{@{}
% %           p{1.2in} rrrrrrrrrrrrr @{}}\hline
% %         \input{/tmp/ahb2.tex}\hline
% %          \end{tabular}\end{scriptsize}\end{table}
% % \end{spacing}



% Table XXX shows variable distributions. If a variable has more than
% 10 categories it is classified into bins...

% %\input .... %TODO !!!! have input here histograms

% \section*{Additional Descriptive Statistics}
% \label{app_des_sta}

% %make sure i have [H] or h! ???
% % \begin{table}[H]
% % \caption{}
% % \centering
% % \label{}
% % \begin{scriptsize}
% % \input{../out/reg_c.tex}
% % \end{scriptsize}
% % \end{table}

%\newpage
%\theendnotes


% \begin{scriptsize}
% \begin{spacing}{.9} 
% \begin{verbatim}
% . tabstat ls, by(lwd7) stat(mean sd) format(%9.1f)

% Summary for variables: ls
% Group variable: lwd7 (RECODE of lwd (do what you like))

%   lwd7 |      Mean        SD
% -------+--------------------
%    0-4 |       5.6       2.7
%      2 |       7.6       2.3
%      3 |       7.6       1.8
%      4 |       8.0       1.5
%      5 |       7.8       1.9
%      6 |       8.2       1.5
%      7 |       8.9       1.5
% -------+--------------------
%  Total |       8.3       1.9
% ----------------------------
% \end{verbatim}
% \end{spacing}{.9}
% \end{scriptsize}


\bibliography{/home/aok/papers/root/tex/ebib.bib,powroty.bib,/home/aok/papers/cityAgency/tex/cityAgency.bib}

\end{spacing}
\end{document}
